% GENERATED, edit the French unit or the translation table
\mxfront{stmg-bank-a}

% -----------------------------------------------------------------------------
% Chapitre 1 -- Suites numériques
% -----------------------------------------------------------------------------
\mxexdef{1}{Recognise an arithmetic sequence by differences}{%
A warehouse manager is testing three models. For every natural number $n$, $n$ is the week number.
\begin{itemize}
\item Number of parcels prepared: $a_n=3(2n+7)-n$.
\item Number of customer complaints: $b_n=2n^2-n+4$.
\item Number of pallets still in stock: $c_n=150-7n$.
\end{itemize}
\begin{enumerate}[label=\alph*.]
\item For each of the three sequences, calculate $u_{n+1}-u_n$ (with $u=a$, $b$ or $c$) and state whether the sequence is arithmetic; if so, give its common difference.
\item For each arithmetic sequence, interpret the common difference in context.
\end{enumerate}
}{%
\begin{enumerate}[label=\alph*.]
\item \textbf{Sequence $(a_n)$.} Expanding: $a_n=6n+21-n=5n+21$. Then $a_{n+1}-a_n=5(n+1)+21-(5n+21)=5$. The difference is constant: $(a_n)$ is arithmetic with common difference $5$ (first term $a_0=21$).

\textbf{Sequence $(b_n)$.} $b_{n+1}-b_n=2(n+1)^2-(n+1)+4-(2n^2-n+4)=2(2n+1)-1=4n+1$. This difference depends on $n$: for example $b_0=4$, $b_1=5$, $b_2=10$, so $b_1-b_0=1$ and $b_2-b_1=5$. The sequence $(b_n)$ is not arithmetic.

\textbf{Sequence $(c_n)$.} $c_{n+1}-c_n=150-7(n+1)-(150-7n)=-7$. The sequence $(c_n)$ is arithmetic with common difference $-7$ (first term $c_0=150$).
\item Each week, the warehouse prepares $5$ more parcels than the week before. Each week, the stock goes down by $7$ pallets (a negative common difference means a decrease).
\end{enumerate}
}{4}

\mxexdef{2}{General term and distant terms (arithmetic)}{%
A delivery company is gradually replacing its diesel vehicles with electric ones. At the start of 2025 (rank $0$), it owns $186$ diesel vehicles and it withdraws $4$ of them every year. Let $d_n$ be the number of diesel vehicles at the start of the year $2025+n$; the sequence $(d_n)$ is arithmetic.
\begin{enumerate}[label=\alph*.]
\item Give the first term $d_0$ and the common difference of the sequence $(d_n)$.
\item Express $d_n$ in terms of $n$ for every natural number $n$.
\item Calculate $d_{12}$ and $d_{30}$. Interpret $d_{30}$ in a sentence.
\end{enumerate}
}{%
\begin{enumerate}[label=\alph*.]
\item $d_0=186$; $4$ vehicles are withdrawn each year, so $d_{n+1}=d_n-4$: the common difference is $r=-4$.
\item For an arithmetic sequence with first term $d_0$: $d_n=d_0+nr$, that is $d_n=186-4n$.
\item $d_{12}=186-4\times12=186-48=138$ and $d_{30}=186-4\times30=186-120=66$.

Rank $30$ corresponds to the year $2025+30=2055$: according to this model, the company will still own $66$ diesel vehicles at the start of $2055$.
\end{enumerate}
}{4}

\mxexdef{3}{Sum of consecutive integers}{%
\begin{enumerate}[label=\alph*.]
\item Calculate $\displaystyle\sum_{i=1}^{365}i=1+2+\cdots+365$.
\item Calculate $\displaystyle\sum_{k=100}^{365}k=100+101+\cdots+365$. State the number of terms in this sum.
\item To build up some savings, Lova takes on a challenge: on the $k$-th day of the year, she puts $k$ cents into a money box ($1$ cent on day 1, $2$ cents on day 2, and so on). How much, in euros, will she have put in altogether over a $365$-day year? And from day 100 to day 365 inclusive?
\end{enumerate}
}{%
\begin{enumerate}[label=\alph*.]
\item The integers from $1$ to $365$ form an arithmetic sequence with common difference $1$; there are $365$ terms. So
$\displaystyle\sum_{i=1}^{365}i=365\times\frac{1+365}{2}=365\times183=\num{66795}$.
\item From $100$ to $365$ there are $365-100+1=266$ terms. So
$\displaystyle\sum_{k=100}^{365}k=266\times\frac{100+365}{2}=266\times\num{232,5}=\num{61845}$.

Check: $\num{66795}-\num{61845}=\num{4950}=\dfrac{99\times100}{2}$, which is indeed the sum of the integers from $1$ to $99$.
\item The total for the year is $\num{66795}$ cents, that is €$\num{667,95}$. From day 100 to day 365, Lova puts in $\num{61845}$ cents, that is €$\num{618,45}$.
\end{enumerate}
}{5}

\mxexdef{4}{General term of a geometric sequence}{%
An app is being tested before its launch. In week $0$, $3$ testers install it. We assume that, each week, the number of new testers is multiplied by $4$. Let $v_n$ be the number of new testers in week $n$; the sequence $(v_n)$ is geometric.
\begin{enumerate}[label=\alph*.]
\item Give the first term and the common ratio of the sequence $(v_n)$, then express $v_n$ in terms of $n$.
\item Calculate $v_5$ and $v_9$.
\item Does the model seem realistic over a long period? Justify in a sentence.
\end{enumerate}
}{%
\begin{enumerate}[label=\alph*.]
\item $v_0=3$ and $v_{n+1}=4v_n$: the common ratio is $q=4$. For a geometric sequence, $v_n=v_0q^n$, so $v_n=3\times4^n$.
\item $v_5=3\times4^5=3\times\num{1024}=\num{3072}$ and $v_9=3\times4^9=3\times\num{262144}=\num{786432}$.
\item No: according to the model there would be $\num{786432}$ new testers in week $9$, and over $3$ million the following week ($v_{10}=\num{3145728}$). Such fast growth cannot last: the number of potential users is limited.
\end{enumerate}
}{4}

\mxexdef{5}{General term and geometric sum}{%
A reforestation association plants $\num{1500}$ trees in the first year of its programme. Thanks to new volunteers, the number of trees planted doubles every year. Let $u_0=\num{1500}$ be the number of trees planted in the first year and $u_n$ the number of trees planted in year $(n+1)$.
\begin{enumerate}[label=\alph*.]
\item Justify that $(u_n)$ is a geometric sequence. Express $u_n$ in terms of $n$ and deduce $u_6$.
\item Calculate $S=u_0+u_1+\cdots+u_6$ and interpret this result in context.
\end{enumerate}
}{%
\begin{enumerate}[label=\alph*.]
\item Doubling means multiplying by $2$: $u_{n+1}=2u_n$. The sequence $(u_n)$ is geometric with first term $u_0=\num{1500}$ and common ratio $q=2$. So $u_n=\num{1500}\times2^n$ and $u_6=\num{1500}\times2^6=\num{1500}\times64=\num{96000}$.
\item The sum has $7$ terms (from $u_0$ to $u_6$) and $q\neq1$:
\[
S=u_0\,\frac{1-q^{7}}{1-q}=\num{1500}\times\frac{1-2^7}{1-2}=\num{1500}\times\frac{-127}{-1}=\num{190500}.
\]
Over the first seven years of the programme, the association planted $\num{190500}$ trees in total.
\end{enumerate}
}{5}

\mxexdef{6}{Choose an arithmetic or geometric model}{%
For each statement, say whether it is true or false. Justify your answer.
\begin{enumerate}
\item The price of an annual travel pass goes up by €$\num{1,50}$ every year. This situation can be modelled by an arithmetic sequence with common difference $\num{1,5}$.
\item A car loses $\pc{15}$ of its value every year. This situation can be modelled by a geometric sequence with common ratio $\num{1,15}$.
\item A sum is invested in a savings account with compound interest at an annual rate of $\pc{3}$, with no deposits or withdrawals. The capital can be modelled by a geometric sequence with common ratio $\num{1,03}$.
\item A rainwater tank holds $\num{1000}$~L. Each day it loses $\pc{5}$ of its contents, then $40$~L are added. The volume can be modelled by a geometric sequence with common ratio $\num{0,95}$.
\item The population of a town grows by $\pc{2}$ per year. This situation can be modelled by an arithmetic sequence with common difference $2$.
\item A loan of €$\num{2000}$ is granted with simple interest at an annual rate of $\pc{6}$. The amount owed can be modelled by an arithmetic sequence with common difference $120$.
\item The number of monthly visits to a website decreases by $300$ every month. This situation can be modelled by a geometric sequence.
\end{enumerate}
}{%
\begin{enumerate}
\item \textbf{True.} €$\num{1,50}$ is added every year: $u_{n+1}=u_n+\num{1,5}$, the absolute change is constant.
\item \textbf{False.} A $\pc{15}$ decrease corresponds to the multiplier $1-\dfrac{15}{100}=\num{0,85}$: the sequence is geometric with common ratio $\num{0,85}$, not $\num{1,15}$ (which would mean an increase).
\item \textbf{True.} Each year the capital is multiplied by $1+\dfrac{3}{100}=\num{1,03}$: $C_{n+1}=\num{1,03}\,C_n$.
\item \textbf{False.} We have $V_{n+1}=\num{0,95}\,V_n+40$: an addition comes on top of the multiplication. With $V_0=\num{1000}$: $V_1=990$ and $V_2=\num{980,5}$; $\dfrac{V_1}{V_0}=\num{0,99}$ and $\dfrac{V_2}{V_1}\approx\num{0,9904}$ are not equal, so the sequence is not geometric (nor is it arithmetic: $-10\neq-\num{9,5}$).
\item \textbf{False.} A $\pc{2}$ rise per year is a constant relative change: the sequence is geometric with common ratio $\num{1,02}$.
\item \textbf{True.} With simple interest, the interest is calculated on the initial amount: $\num{2000}\times\num{0,06}=120$ euros are added each year. The amount owed therefore goes up by $120$ euros per year.
\item \textbf{False.} The decrease is absolute and constant ($-300$ per month): the sequence is arithmetic with common difference $-300$.
\end{enumerate}
}{6}

% -----------------------------------------------------------------------------
% Chapitre 2 -- Fonctions exponentielles
% -----------------------------------------------------------------------------
\mxexdef{7}{Extend a geometric sequence to a function}{%
On 1 January 2024, the capacity of the solar panels installed in a region is $240$ megawatts (MW). This capacity is estimated to grow by $\pc{12}$ per year. For every natural number $n$, let $u_n$ be the installed capacity, in MW, on 1 January of the year $2024+n$; we assume that $(u_n)$ is a geometric sequence.
\begin{enumerate}
\item Give the common ratio of the sequence $(u_n)$ and express $u_n$ in terms of $n$.
\item We want to estimate the installed capacity at any date. Suggest an expression for a function $f$ giving the installed capacity $x$ years after 1 January 2024, where $x$ is a positive real number.
\item Estimate, to the nearest $\num{0,1}$~MW, the installed capacity on 1 January 2026, on 1 April 2027 and on 1 July 2029. A quarter is taken as $\num{0,25}$ year.
\end{enumerate}
}{%
\begin{enumerate}
\item A $\pc{12}$ increase corresponds to the multiplier $1+\dfrac{12}{100}=\num{1,12}$. The common ratio is $q=\num{1,12}$ and $u_n=240\times\num{1,12}^n$.
\item We extend the sequence to positive real numbers: $f(x)=240\times\num{1,12}^x$ (indeed $f(n)=u_n$ for every integer $n$).
\item \begin{itemize}
\item 1 January 2026: $x=2$ and $f(2)=240\times\num{1,12}^2=\num{301,056}\approx\num{301,1}$~MW.
\item 1 April 2027: $3$ years and one quarter, so $x=\num{3,25}$; $f(\num{3,25})=240\times\num{1,12}^{\num{3,25}}\approx\num{346,9}$~MW.
\item 1 July 2029: $5$ and a half years, so $x=\num{5,5}$; $f(\num{5,5})=240\times\num{1,12}^{\num{5,5}}\approx\num{447,6}$~MW.
\end{itemize}
According to this model, the installed capacity would be about $\num{301,1}$~MW, $\num{346,9}$~MW and $\num{447,6}$~MW on these three dates.
\end{enumerate}
}{5}

\mxexdef{8}{Algebraic rules for powers}{%
Write each of the following numbers as a single power of $3$:
\[
a=3^{-2}\times3^{6};\qquad b=\left(3^{\num{1,5}}\right)^4;\qquad
c=27\times9^{-2};\qquad d=\frac{3^8}{9^3}.
\]
}{%
\begin{itemize}
\item $a=3^{-2+6}=3^4$ (rule $a^xa^y=a^{x+y}$).
\item $b=3^{\num{1,5}\times4}=3^6$ (rule $(a^x)^y=a^{xy}$).
\item $27=3^3$ and $9=3^2$, so $9^{-2}=\left(3^2\right)^{-2}=3^{-4}$; hence $c=3^3\times3^{-4}=3^{-1}$.
\item $9^3=\left(3^2\right)^3=3^6$, so $d=\dfrac{3^8}{3^6}=3^{8-6}=3^2$.
\end{itemize}
This can be checked on a calculator, for example $c=\dfrac{27}{81}=\dfrac13=3^{-1}$ and $d=\dfrac{\num{6561}}{729}=9=3^2$.
}{3}

\mxexdef{9}{Direction of variation of $t\mapsto ka^t$}{%
Determine the direction of variation on $\mathbb{R}$ of the functions defined by:
\[
f:t\mapsto\num{3,5}\times\num{0,82}^t;\qquad
g:t\mapsto-\num{0,6}\times\num{1,15}^t;\qquad
h:t\mapsto-4\times\num{0,7}^t.
\]
}{%
\begin{itemize}
\item $0<\num{0,82}<1$, so $t\mapsto\num{0,82}^t$ is strictly decreasing; the coefficient $\num{3,5}$ is positive and keeps the direction of variation: $f$ is strictly decreasing.
\item $\num{1,15}>1$, so $t\mapsto\num{1,15}^t$ is strictly increasing; the coefficient $-\num{0,6}$ is negative and reverses the direction of variation: $g$ is strictly decreasing.
\item $0<\num{0,7}<1$, so $t\mapsto\num{0,7}^t$ is strictly decreasing; the coefficient $-4$ is negative and reverses the direction of variation: $h$ is strictly increasing.
\end{itemize}
Numerical check: $h(0)=-4$, $h(1)=-\num{2,8}$ and $h(2)=-\num{1,96}$: the values increase.
}{3}

\mxexdef{10}{Average annual rate of change}{%
Over five years, a factory has reduced its water consumption by $\pc{27}$.
\begin{enumerate}[label=\alph*.]
\item Give the overall multiplier over the five years.
\item Calculate the average annual rate of change of the water consumption, to the nearest $\pc{0,1}$.
\item Check the result by applying this rounded average multiplier five times.
\end{enumerate}
}{%
\begin{enumerate}[label=\alph*.]
\item A $\pc{27}$ decrease corresponds to the overall multiplier $Q=1-\dfrac{27}{100}=\num{0,73}$.
\item If $q$ is the average annual multiplier, then $q^5=\num{0,73}$, so $q=\num{0,73}^{\frac15}\approx\num{0,9390}$. The average annual rate is $t=q-1\approx-\num{0,061}$: water consumption fell on average by about $\pc{6,1}$ per year.
\item $\num{0,939}^5\approx\num{0,7300}$, which does correspond to an overall decrease of about $\pc{27}$.
\end{enumerate}
}{4}

\mxexdef{11}{Reading graphs of $a^x$ curves}{%
The graph below shows, in a coordinate system, the functions $f$ and $g$ defined on $\mathbb{R}$ by $f(x)=\num{1,8}^x$ and $g(x)=\num{0,45}^x$, together with the line with equation $y=3$.
\begin{center}
\begin{tikzpicture}
\begin{axis}[width=.78\textwidth,height=6.5cm,axis lines=middle,xmin=-3,xmax=3,ymin=0,ymax=6,
  xtick={-3,...,3},ytick={0,...,6},minor x tick num=4,minor y tick num=1,
  grid=both,major grid style={black!25},minor grid style={black!10},xlabel={$x$},ylabel={$y$}]
\addplot[mxteacher,very thick,samples=120,domain=-3:3]{0.45^x};
\addplot[mxwarn,very thick,samples=120,domain=-3:3]{1.8^x};
\addplot[mxmint,thick,dashed,domain=-3:3]{3};
\node[mxteacher,font=\small] at (axis cs:-1.45,4.8) {$\mathcal{C}_1$};
\node[mxwarn,font=\small] at (axis cs:2.2,5.2) {$\mathcal{C}_2$};
\node[mxmint,font=\small,fill=white,inner sep=1pt] at (axis cs:0.5,3.3) {$y=3$};
\end{axis}
\end{tikzpicture}
\end{center}
\begin{enumerate}
\item Which curve represents the function $f$? Justify your choice.
\item Solve the equation $f(x)=3$ graphically. Using a table of values on the calculator, give an approximate value of the solution to the nearest $\num{0,01}$.
\item Determine graphically the set of real numbers $x$ such that $f(x)\geqslant3$, then the set of real numbers $x$ such that $g(x)\geqslant3$.
\end{enumerate}
}{%
\begin{enumerate}
\item $\num{1,8}>1$, so $f$ is strictly increasing; $0<\num{0,45}<1$, so $g$ is strictly decreasing. The increasing curve $\mathcal{C}_2$ represents $f$ (and $\mathcal{C}_1$ represents $g$). Both curves pass through the point $\pt{0;1}$.
\item The curve $\mathcal{C}_2$ meets the line $y=3$ at a point whose \textit{x}-coordinate reads close to $\num{1,9}$. With the calculator: $\num{1,8}^{\num{1,86}}\approx\num{2,984}$ and $\num{1,8}^{\num{1,87}}\approx\num{3,002}$, so the solution lies between $\num{1,86}$ and $\num{1,87}$; the value $\num{1,87}$ gives the image closest to $3$. The solution is $x\approx\num{1,87}$.
\item $f(x)\geqslant3$ when $\mathcal{C}_2$ is above the line $y=3$, that is for $x\geqslant\num{1,87}$ approximately: $\intfo{\num{1,87}}{+\infty}$ (approximate value of the endpoint).

$g$ is decreasing and $\mathcal{C}_1$ meets the line $y=3$ at an abscissa close to $-\num{1,4}$ (calculator: $\num{0,45}^{-\num{1,38}}\approx\num{3,010}$ and $\num{0,45}^{-\num{1,37}}\approx\num{2,986}$, so $x\approx-\num{1,38}$). Hence $g(x)\geqslant3$ for $x\leqslant-\num{1,38}$ approximately: $\intof{-\infty}{-\num{1,38}}$ (approximate value of the endpoint).
\end{enumerate}
}{5}

\mxexdef{12}{Inequalities with the same base}{%
Solve the following inequalities on $\mathbb{R}$.
\begin{enumerate}[label=\alph*.]
\item $\num{0,7}^x\leqslant\num{0,7}^{-1}$
\item $\num{2,5}^{x+1}>\num{2,5}^4$
\item $\num{0,5}^x\geqslant\num{0,25}$
\end{enumerate}
}{%
\begin{enumerate}[label=\alph*.]
\item $0<\num{0,7}<1$: the function $x\mapsto\num{0,7}^x$ is strictly decreasing, so it reverses the order. Hence $\num{0,7}^x\leqslant\num{0,7}^{-1}\iff x\geqslant-1$. Solution set: $\intfo{-1}{+\infty}$.
\item $\num{2,5}>1$: the function $x\mapsto\num{2,5}^x$ is strictly increasing, so it preserves the order. Hence $\num{2,5}^{x+1}>\num{2,5}^4\iff x+1>4\iff x>3$. Solution set: $\intoo{3}{+\infty}$.
\item $\num{0,25}=\num{0,5}^2$, so the inequality reads $\num{0,5}^x\geqslant\num{0,5}^2$. Since $x\mapsto\num{0,5}^x$ is strictly decreasing: $x\leqslant2$. Solution set: $\intof{-\infty}{2}$.
\end{enumerate}
}{3}

\mxexdef{13}{Model with $f(t)=ka^t$ and compute}{%
Just after an injection, the concentration of a drug in a patient's blood is $80$~mg/L. This concentration then decreases by $\pc{15}$ every hour. This is an exponential decay.

It is modelled by a function $C$ of the time $t$, in hours, such that $C(t)=ka^t$.
\begin{enumerate}
\item Determine the constants $k$ and $a$.
\item Calculate the concentration after $3$~h~$30$~min, to the nearest $\num{0,1}$~mg/L.
\item The doctor plans a new injection as soon as the concentration drops below $50$~mg/L. Is this already the case after $3$~h~$30$~min? Answer in a sentence.
\end{enumerate}
}{%
\begin{enumerate}
\item $C(0)=k\times a^0=k$, and $C(0)=80$, so $k=80$. A $\pc{15}$ decrease per hour corresponds to the multiplier $1-\dfrac{15}{100}=\num{0,85}$, so $a=\num{0,85}$ and $C(t)=80\times\num{0,85}^t$.
\item $3$~h~$30$~min $=\num{3,5}$~h. $C(\num{3,5})=80\times\num{0,85}^{\num{3,5}}\approx\num{45,3}$~mg/L.
\item $\num{45,3}<50$: after $3$~h~$30$~min the concentration has already dropped below $50$~mg/L, so the new injection is due (in fact $C(3)\approx\num{49,1}$~mg/L, so the concentration is below this threshold from $3$~h onwards).
\end{enumerate}
}{4}

\mxexdef{14}{Average rate of two successive changes}{%
\begin{enumerate}
\item The price of a digital currency rises by $\pc{125}$ one year, then falls by $\pc{36}$ the following year. What is the average annual rate of change over these two years?
\item The price of a raw material rises by $\pc{80}$ one month, then falls by $\pc{55}$ the following month. What is the average monthly rate of change?
\item A student claims that, in question 2, the average rate is $\dfrac{80-55}{2}=\pc{12,5}$. Explain in a sentence why this is wrong.
\end{enumerate}
}{%
\begin{enumerate}
\item The multipliers are $1+\dfrac{125}{100}=\num{2,25}$ and $1-\dfrac{36}{100}=\num{0,64}$. The overall multiplier is $\num{2,25}\times\num{0,64}=\num{1,44}$. The average annual multiplier $q$ satisfies $q^2=\num{1,44}$, so $q=\sqrt{\num{1,44}}=\num{1,2}$: an average annual rise of $\pc{20}$.
\item The multipliers are $\num{1,8}$ and $\num{0,45}$; the overall multiplier is $\num{1,8}\times\num{0,45}=\num{0,81}$. The average multiplier is $\sqrt{\num{0,81}}=\num{0,9}$, that is an average monthly fall of $\pc{10}$.
\item Successive changes combine by multiplying the multipliers, not by adding the rates: the $\pc{55}$ fall applies to a price that has already risen. Over the two months the price is multiplied by $\num{0,81}$ (an overall fall of $\pc{19}$), which corresponds to an average fall of $\pc{10}$ per month, not a rise of $\pc{12,5}$.
\end{enumerate}
}{5}
